\section{Safety, Ethics, and Compliance}

The project contains a computational component and a proposed experimental validation component. The computational component can proceed independently. The physical experiment will not begin until appropriate school approval, local scientific review, risk assessment, and qualified adult supervision are secured.

The proposed experimental validation may involve high voltage and flammable solvent, so it must be treated as a hazardous-device and hazardous-chemical activity. Before purchasing, assembling, or operating any apparatus, the project team will prepare required approval forms, a full risk assessment, safety data sheet documentation, supervision plan, schematic review, operating boundaries, and emergency shutoff procedures. Operation will occur only in a supervised school or laboratory environment, with appropriate enclosure, interlock, current limiting, grounding, ventilation, PPE, and waste-disposal procedures as required by school and competition rules.

No live high-voltage demonstration is planned for a public fair setting. Experimental data, images, videos, and analysis results will be presented instead. The apparatus is intended for microfluidic Taylor-cone observation and image-based validation, not as a projectile device or propulsion demonstration. The project will comply with all applicable Hong Kong selection-fair and Regeneron ISEF-affiliated rules.
